% !TeX root = main.tex

\section{Targeted Inference with Nuisance Parameters}
\label{sec:nuisance}

We extend the sequential contrastive estimation framework to handle nuisance parameters—parameters that must be modeled but are not the focus of inference.

\subsection{Notation}

Partition parameters as $\theta = (\theta_T, \theta_N)$ where $\theta_T$ are targets of interest and $\theta_N$ are nuisance parameters. The prior is $p(\theta) = p(\theta_T, \theta_N)$ with conditional $p(\theta_N|\theta_T)$.
We use subscripts $T,N$ to denote target and nuisance blocks, and superscripts $(\cdot)$ to index Monte Carlo samples.

The \emph{targeted expected information gain} is:
\begin{equation}
\mathcal{I}_K^{\text{tgt}}(\pi) = I(\theta_T; h_K | \pi) = \E_{p(\theta)p(h_K|\theta,\pi)}\left[\log \frac{p(h_K|\theta_T,\pi)}{p(h_K|\pi)}\right]
\end{equation}
where the marginal likelihood is $p(h_K|\theta_T,\pi) = \int p(h_K|\theta_T, \theta_N, \pi) p(\theta_N|\theta_T) \, d\theta_N$.

\subsection{Practical Nested Monte Carlo Objective}

Since $p(h_K|\theta_T,\pi)$ is intractable, we use Monte Carlo estimation with:
\begin{itemize}
    \item \textbf{Numerator}: $M$ samples $\theta_N^{(1:M)} \sim p(\theta_N|\theta_T^{(0)})$ to estimate $p(h_K|\theta_T^{(0)},\pi)$
    \item \textbf{Denominator}: $L$ contrastive samples $\theta^{(1)},\ldots,\theta^{(L)} \sim p(\theta)$ (full parameter vectors) to estimate $p(h_K|\pi)$
\end{itemize}

\begin{theorem}[Targeted sPCE Lower Bound]
\label{thm:targeted_spce}
Let $\theta^{(0)} = (\theta_T^{(0)}, \theta_N^{(0)}) \sim p(\theta)$ generate $h_K \sim p(h_K|\theta^{(0)}, \pi)$. Let $\theta^{(1)},\ldots,\theta^{(L)} \sim p(\theta)$ be independent contrastive samples, and let $\theta_N^{(1:M)} \sim p(\theta_N|\theta_T^{(0)})$ be nuisance samples. Define
\begin{equation}
\hat{\mathcal{L}}_K^{\text{tgt}}(\pi, L, M) = \E\left[\log \frac{\frac{1}{M}\sum_{m=1}^{M} p(h_K|\theta_T^{(0)}, \theta_N^{(m)}, \pi)}{\frac{1}{L+1}\sum_{\ell=0}^{L} p(h_K|\theta^{(\ell)}, \pi)}\right].
\end{equation}
Then $\hat{\mathcal{L}}_K^{\text{tgt}}(\pi, L, M) \leq \mathcal{I}_K^{\text{tgt}}(\pi)$, with $\hat{\mathcal{L}}_K^{\text{tgt}} \to \mathcal{I}_K^{\text{tgt}}$ as $L, M \to \infty$.
\end{theorem}

\begin{proof}
Define:
\begin{align}
A &= \frac{1}{M}\sum_{m=1}^{M} p(h_K|\theta_T^{(0)}, \theta_N^{(m)}, \pi) & &\text{(estimator of } p(h_K|\theta_T^{(0)},\pi)\text{)} \\
B &= \frac{1}{L+1}\sum_{\ell=0}^{L} p(h_K|\theta^{(\ell)}, \pi) & &\text{(estimator of } p(h_K|\pi)\text{, includes } \theta^{(0)}\text{)}
\end{align}

The gap between the true objective and the practical bound is:
\begin{align}
\mathcal{I}_K^{\text{tgt}} - \hat{\mathcal{L}}_K^{\text{tgt}} &= \E\left[\log \frac{p(h_K|\theta_T^{(0)},\pi)}{p(h_K|\pi)} - \log \frac{A}{B}\right] \\
&= \E\left[\log \frac{p(h_K|\theta_T^{(0)},\pi)}{A}\right] + \E\left[\log \frac{B}{p(h_K|\pi)}\right]
\end{align}

\textbf{Term 1 (Nuisance marginalization gap):} Condition on the realized $(\theta_T^{(0)}, h_K)$. Since $\theta_N^{(1:M)} \overset{\text{i.i.d.}}{\sim} p(\theta_N|\theta_T^{(0)})$,
\begin{equation}
\E[A \mid \theta_T^{(0)}, h_K] = \int p(h_K|\theta_T^{(0)},\theta_N,\pi)\,p(\theta_N|\theta_T^{(0)})\,d\theta_N = p(h_K|\theta_T^{(0)},\pi).
\end{equation}
Since $\log$ is concave, Jensen's inequality gives
\begin{equation}
\E[\log A \mid \theta_T^{(0)}, h_K] \leq \log \E[A \mid \theta_T^{(0)}, h_K] = \log p(h_K|\theta_T^{(0)},\pi),
\end{equation}
and hence $\E[\log(p(h_K|\theta_T^{(0)},\pi)/A) \mid \theta_T^{(0)}, h_K] \geq 0$. Taking expectations over $(\theta_T^{(0)}, h_K)$ yields $\E[\log(p(h_K|\theta_T^{(0)},\pi)/A)] \geq 0$.

\textbf{Term 2 (Standard sPCE gap):} This is exactly the gap analyzed in Foster et al.~\cite{foster2021}. Using the symmetry of $\theta^{(0)}$ being included in the denominator sum and a change of measure argument, this equals:
\begin{equation}
\E_{p(h_K|\pi)}\left[\text{KL}\left(\tilde{p}(\theta^{(0:L)}|h_K) \,\|\, p(\theta^{(0:L)})\right)\right] \geq 0
\end{equation}
where $\tilde{p}$ is a mixture distribution (see~\cite{foster2021} for details).

Since both terms are non-negative: $\hat{\mathcal{L}}_K^{\text{tgt}} \leq \mathcal{I}_K^{\text{tgt}}$.

\textbf{Convergence:} Assume there exist constants $0<\kappa_{1,N} \le \kappa_{2,N}<\infty$ and $0<\kappa_1 \le \kappa_2<\infty$ such that
\begin{align}
\kappa_{1,N} \le \frac{p(h_K|\theta_T,\theta_N,\pi)}{p(h_K|\theta_T,\pi)} \le \kappa_{2,N} && \forall\,\theta_T,\theta_N,h_K, \\
\kappa_1 \le \frac{p(h_K|\theta,\pi)}{p(h_K|\pi)} \le \kappa_2 && \forall\,\theta,h_K,
\end{align}
as in the standard sPCE analysis~\cite{foster2021}. Then the log-ratios in Term 1 and Term 2 are bounded. By the Strong Law of Large Numbers, $A \to p(h_K|\theta_T^{(0)},\pi)$ and $B \to p(h_K|\pi)$ almost surely as $M,L\to\infty$, and the Bounded Convergence Theorem yields Term 1 $\to 0$ and Term 2 $\to 0$; moreover Term 2 decreases at rate $\mathcal{O}(L^{-1})$.
\end{proof}

\begin{remark}
The practical objective requires evaluating $p(h_K|\theta, \pi)$ for $(1+L) + M$ parameter settings per training iteration: $L+1$ full samples for the denominator, plus $M$ nuisance samples for the numerator (all sharing the same $\theta_T^{(0)}$).
\end{remark}

\subsection{Computational Budget Trade-offs}
\label{sec:budget_stub}
This appendix section is reserved for a discussion of how to allocate a fixed computational budget across the number of contrastive samples $L$, nuisance samples $M$, and minibatch size $B$.
In particular, it will summarize how the bound tightness and Monte Carlo gaps scale with $L$ and $M$, and how the variance of stochastic gradient estimates scales with $B$.
